\section{Uniform、Beta 与 Dirichlet 分布}
\label{sec:05-Probability/06-Common-Distributions/04}

推荐位 A 展示了 \(10\) 次，被点了 \(7\) 次，点击率 \(0.7\)。推荐位 B 展示了 \(10000\) 次，被点了 \(6000\) 次，点击率 \(0.6\)。按点击率排序，A 排在前面。

没人会真的这么排。\(0.7\) 这个数背后只有十次曝光，换一天可能就是 \(0.4\)；\(0.6\) 背后有一万次曝光，明天大概还是 \(0.6\)。两个比值的可信程度差着一个数量级，而``比值''这个形式把这个差别完全丢掉了。

要把它捡回来，需要的不是一个更好的点估计，而是一个关于 \(p\) 本身的分布——一个定义在 \([0,1]\) 上、能随着数据变窄的对象。本节讲的就是这类分布，以及把它们推广到多个类别之后的样子。

\subsection{从一条平线开始，以及``均匀''的歧义}

最省事的起点是均匀分布：\(p\) 在 \([0,1]\) 上处处等可能，密度恒为 \(1\)。区间概率只由长度决定，均值 \(1/2\)，方差 \(1/12\)。作为``什么都不知道''的表达，它看上去无可指摘。

但``均匀''这个词依赖于你用什么刻度。如果 \(X\) 在 \([1,10]\) 上均匀，\(\log X\) 就不均匀；反过来也一样。落到概率参数上：说``对 \(p\) 一无所知''时，你是指 \(p\) 均匀，还是几率 \(p/(1-p)\) 均匀，还是对数几率 \(\log\bigl(p/(1-p)\bigr)\) 均匀？三种说法各自对应一个不同的先验，第一个是 \(\operatorname{Beta}(1,1)\)，最后一个对应 \(\operatorname{Beta}(0,0)\)——一个在两端发散、不能正常使用的极限对象。追求参数化不变性的 Jeffreys 先验给出 \(\operatorname{Beta}(0.5,0.5)\)，两端同样发散但慢得多。三者都自称不偏不倚，含义却各不相同。

所以``取均匀先验''不是一个免于选择的选项，它只是把选择藏进了坐标的挑选里。承认这一点之后，与其纠结哪条平线最中立，不如直接用一个形状可调的族，把``我大概相信什么''写成参数。

\subsection{Beta：比例的分布}

Beta 分布的密度是

\[
f(p)=\frac{1}{B(\alpha,\beta)}\,p^{\alpha-1}(1-p)^{\beta-1},\qquad 0<p<1,
\]

其中 \(B(\alpha,\beta)=\Gamma(\alpha)\Gamma(\beta)/\Gamma(\alpha+\beta)\) 只负责把总质量归到 \(1\)，不影响形状。形状全在 \(p^{\alpha-1}(1-p)^{\beta-1}\) 这个核里：\(\alpha\) 把质量往 \(1\) 那头拉，\(\beta\) 往 \(0\) 那头拉。两个都等于 \(1\) 时核变成常数，退回均匀分布；两个都大于 \(1\) 时质量集中在内部，形如一座小山；任一个小于 \(1\) 时对应那一端的密度发散，表达``我认为 \(p\) 很可能贴着边界''。

均值和方差写成这两个形式更好用：

\[
\E[P]=\frac{\alpha}{\alpha+\beta},\qquad
\Var(P)=\frac{\alpha\beta}{(\alpha+\beta)^2(\alpha+\beta+1)}.
\]

比值 \(\alpha/(\alpha+\beta)\) 决定中心落在哪，总量 \(\alpha+\beta\)（常称浓度）决定分布围着中心有多紧。方向和确信度被分给两个参数，各调各的——这正是均匀分布做不到的事。

\begin{center}
\begin{tikzpicture}[x=1cm,y=1cm,font=\footnotesize]
  \draw[black!70,->] (0,0) -- (8.6,0);
  \draw[black!70,->] (0,0) -- (0,4.2);
  \foreach \p/\s in {0/0, 0.2/0.2, 0.4/0.4, 0.6/0.6, 0.8/0.8, 1/1}{
     \draw[black!60] ({8*\p},0) -- ({8*\p},-0.1);
     \node[below,black!60,font=\scriptsize] at ({8*\p},-0.1) {\s};}
  \node[below,black!70] at (4,-0.55) {点击率 \(p\)};
  \draw[black!55,dashed,thick] (0,1.2) -- (8,1.2);
  \node[black!60] at (1.5,1.62) {先验 \(\operatorname{Beta}(1,1)\)};
  \draw[blue!55!black,thick,domain=0:1,samples=90]
     plot ({8*\x},{1.2*60*\x*\x*\x*(1-\x)*(1-\x)});
  \node[blue!55!black] at (2.6,2.95) {\(3\) 点 \(2\) 未点：\(\operatorname{Beta}(4,3)\)};
  \draw[red!65!black,thick,domain=0:1,samples=90]
     plot ({8*\x},{1.2*1320*\x*\x*\x*\x*\x*\x*\x*(1-\x)*(1-\x)*(1-\x)});
  \node[red!65!black] at (6.4,3.9) {\(7\) 点 \(3\) 未点：\(\operatorname{Beta}(8,4)\)};
\end{tikzpicture}

\emph{同一个推荐位，观测一条条进来。起点是平的先验；五次曝光之后隆起一个宽包，中心在 \(0.6\)；十次曝光之后中心移到 \(0.7\)，包也窄了一些。注意十次曝光换来的仍是一个很宽的分布——图上 \(0.45\) 到 \(0.9\) 之间都有可观的密度，所以那个 \(0.7\) 完全撑不起``排第一''的结论。}
\end{center}

\subsection{更新之后还在同一个族里}

上图的三条曲线之间只差一次乘法。设 \(p\) 的先验是 \(\operatorname{Beta}(\alpha,\beta)\)，观测到 \(s\) 次成功 \(f\) 次失败，似然中与 \(p\) 有关的部分是 \(p^s(1-p)^f\)。乘上先验的核：

\[
p^{\alpha-1}(1-p)^{\beta-1}\cdot p^{s}(1-p)^{f}=p^{\alpha+s-1}(1-p)^{\beta+f-1}.
\]

右边正是 \(\operatorname{Beta}(\alpha+s,\beta+f)\) 的核。后验没有跑出 Beta 族，更新规则简单到不需要计算：成功数加给 \(\alpha\)，失败数加给 \(\beta\)。这种性质叫共轭（关于它在更一般模型里的样子，见\hyperref[sec:13-Machine-Learning/06-Probabilistic-Models/07]{``共轭先验让 Bayesian 更新保持同一形式''}）。

由此可以把 \(\alpha\) 和 \(\beta\) 读成伪计数：动手采集数据之前，你的信念相当于已经虚拟地看过 \(\alpha-1\) 次成功和 \(\beta-1\) 次失败。参数不是整数时这只是个记忆装置，但方向是对的——先验的分量就是用``它相当于几条观测''来衡量的。

后验均值把这层意思写得更直白：

\[
\E[p\given \text{数据}]=\frac{\alpha+s}{\alpha+\beta+s+f}.
\]

它是先验均值和样本比例之间的加权平均，权重就是伪计数与真实计数之比。数据少时先验说了算，数据多了样本比例接管。这种收缩恰好挡住了小样本最典型的失误：一次曝光一次点击就报 \(p=1.0\)。

回到那两个推荐位。取一个温和的先验 \(\operatorname{Beta}(2,2)\)（相当于事先见过一次点击一次未点击），A 位的后验是 \(\operatorname{Beta}(9,5)\)，均值 \(0.643\)；B 位的后验是 \(\operatorname{Beta}(6002,4002)\)，均值 \(0.600\)。均值上 A 仍然领先，但两者的后验标准差差了两个数量级：A 约 \(0.12\)，B 约 \(0.005\)。把这个不确定性算进去再排序，结论就完全不同了——线上系统里常用的 Thompson 采样，做的正是从每个候选的后验里各抽一个样本再比大小，于是曝光不足的候选会时不时被抽到高分，从而获得继续被试的机会。共轭在这里的价值是工程性的：后验参数只是两个计数器，可以在流式更新中一直加下去，不需要重新拟合。

\subsection{Dirichlet：同一套逻辑，多个类别}

把二分类换成 \(K\) 类，概率向量 \(p=(p_1,\ldots,p_K)\) 满足 \(p_i\ge0\) 且 \(\sum_i p_i=1\)，住在 \((K-1)\) 维单纯形上。Dirichlet 分布的核就是 Beta 核的直接推广，

\[
\prod_{i=1}^K p_i^{\alpha_i-1},
\]

每个分量配一个伪计数 \(\alpha_i\)。记 \(\alpha_0=\sum_j\alpha_j\)，则 \(\E[p_i]=\alpha_i/\alpha_0\)，浓度仍由 \(\alpha_0\) 承担。对 Multinomial 计数 \(n_1,\ldots,n_K\)，更新方式和 Beta 完全一样：后验是 \(\operatorname{Dirichlet}(\alpha_1+n_1,\ldots,\alpha_K+n_K)\)。\(K=2\) 时它就是 Beta。

分量之间必然负相关，这不是参数选出来的，是单纯形的几何逼出来的：总和锁死为 \(1\)，一个分量涨，别的只能让位。相应地，相关结构完全没有自由度——任意两个分量的相关系数被同一组 \(\alpha\) 锁定成同一个函数形式，你无法表达``第一类和第二类强烈互斥，第三类基本独立''这种模式。

浓度参数的取值在实践中会明显改变结果的样子。\(\alpha_i\) 全部大于 \(1\) 时抽出的概率向量偏向均衡，各类都分到一点；\(\alpha_i\) 全部小于 \(1\) 时抽出的向量趋于稀疏，质量集中在少数几个分量上。主题模型里的文档--主题先验通常取 \(\alpha_i\) 远小于 \(1\)，为的正是这个效果——一篇文档只讲两三个主题，而不是把所有主题平均掺一点。另一头，自然语言处理里的加一平滑其实就是 \(\operatorname{Dirichlet}(1,\ldots,1)\) 先验下的后验均值：没见过的词不再被判零概率，因为先验替它预付了一次计数。

Dirichlet 还有一个构造值得记住：取 \(K\) 个独立的 \(\operatorname{Gamma}(\alpha_i,1)\)，除以它们的总和，得到的向量精确服从 \(\operatorname{Dirichlet}(\alpha_1,\ldots,\alpha_K)\)。这既给了一个现成的采样方法，也给了一个物理解读——每个类别先各自拥有一份绝对质量，归一化只是把绝对量换算成比例，而总质量越大，换算出的比例越稳定。上一节的 Gamma 在这里第二次出场，这次它扮演的是零件而不是等待时间。

\subsection{共轭买到了什么，又押上了什么}

共轭买到的是闭式：后验永远能写出来，永远只是加几个计数，更新可以流式进行，超参数有可解释的含义。这在需要对成千上万个候选各维护一个后验的场合（每个广告位、每个商品、每个臂）几乎不可替代。

押上的是形状。Beta 只能是单峰或两端发散的那几种样子，Dirichlet 只有一种相关结构。如果真实的后验应该有两个峰——比如某个转化率要么很低要么很高，中间不太可能——单个 Beta 表达不了，需要换成混合模型。如果你需要任意协方差结构，常见的替代是对概率向量取对数比值之后假设多元 Gaussian（Logistic-normal），它换来了灵活性，要付出的是后验不再有闭式，得靠变分或采样。

还有一条更基本的提醒。\(\operatorname{Dirichlet}(1,\ldots,1)\) 在单纯形上确实是均匀的，但正如本节开头所说，均匀是相对于坐标而言的：你关心的如果是``最大分量有多大''或者``前三类的概率之和''，这些导出量在这个先验下并不均匀。用它是因为你真的相信所有概率向量先验等可能，还是因为它计算方便——两个理由通向不同的分析，值得对自己说清楚。

到这里为止，本章遇到的分布都被支撑集牢牢框住：整数、非负实数、\([0,1]\) 区间、单纯形。下一节的分布铺满整条实轴，没有任何硬边界。它之所以成为默认选项，靠的是两个彼此独立的理由，而这两个理由各自都有失效的条件。
